Следующий шаг проверял ещё более узкую гипотезу: является ли режим
`group\_rows=512` уже практическим потолком, или дальнейшее увеличение числа
строк в группе всё ещё уменьшает latency быстрее, чем растёт локальный union
значений. Для этого был выполнен дополнительный frontier-sweep по
`group\_rows=256,512,1024,2048` и `group\_cols=4096,8192`, при котором,
помимо времени, измерялись также средний и суммарный local palette union.

Результат оказался ещё сильнее: `group\_rows > 512` продолжает заметно
помогать. Для двух approved слоёв характерные точки выглядят так:

\begin{center}
\begin{tabular}{llrrrrr}
\toprule
Layer & $(group\_rows, group\_cols)$ & launches & total union & grouped-global ms & grouped-local ms & output MSE \\
\midrule
q\_proj & $(256, 8192)$ & 32 & 46828 & 9.70 & 8.74 & $4.58 \times 10^{-6}$ \\
q\_proj & $(512, 8192)$ & 16 & 24791 & 5.67 & 5.15 & $4.58 \times 10^{-6}$ \\
q\_proj & $(1024, 8192)$ & 8 & 12658 & 3.30 & 3.17 & $4.58 \times 10^{-6}$ \\
q\_proj & $(2048, 8192)$ & 4 & 6370 & 1.90 & 1.82 & $4.58 \times 10^{-6}$ \\
up\_proj & $(256, 8192)$ & 112 & 140824 & 29.82 & 26.63 & $1.36 \times 10^{-6}$ \\
up\_proj & $(512, 8192)$ & 56 & 78579 & 17.05 & 16.11 & $1.36 \times 10^{-6}$ \\
up\_proj & $(1024, 8192)$ & 28 & 41714 & 11.61 & 11.09 & $1.36 \times 10^{-6}$ \\
up\_proj & $(2048, 8192)$ & 14 & 21315 & 6.82 & 6.44 & $1.36 \times 10^{-6}$ \\
\bottomrule
\end{tabular}
\end{center}

Эти числа важны сразу в двух отношениях. Во-первых, лучший pure-latency режим
теперь смещается к `(2048,8192)`: local path достигает `1.82` мс на `q_proj`
и `6.44` мс на `up_proj`, то есть остаётся лишь примерно в `1.27x-1.30x`
медленнее, чем fused global full Triton. Во-вторых, статистика union меняет
саму интерпретацию bottleneck-а: средний union на одну группу растёт лишь
умеренно, тогда как суммарный union по всему слою резко падает, потому что
число launches сокращается быстрее. Следовательно, оставшийся runtime-gap уже
хуже объясняется ``слишком большим local union'' как таковым и лучше
объясняется эффективностью арифметики и scheduling quality grouped local
kernel на очень больших 2D blocks.

Если же нужен более консервативный near-frontier режим, который уже очень близок
к лучшему latency, но не требует максимально агрессивного row grouping, то таким
режимом выглядит `(1024,8192)`: `3.17` мс на `q_proj` и `11.09` мс на
`up_proj`, всё ещё с небольшим преимуществом local path над grouped-global.

После этого была проверена ещё одна очень узкая systems-гипотеза. Сначала
autotune-space local Triton kernel был расширен более широкими block-конфигурациями,
но это практически ничего не изменило: frontier latency сдвинулась менее чем на
`0.1\%`. Значит, bottleneck уже не объяснялся банально отсутствующим autotune
config.

Следующий ближайший кандидат оказался важнее: local palette indices по-прежнему
хранились как `int32`, хотя измеренные local unions на практике были значительно
меньше предела `int16`. После перехода на `int16` для local indices были
получены следующие изменения:

\begin{center}
\begin{tabular}{llrrr}
\toprule
Layer & $(group\_rows, group\_cols)$ & old local ms & new local ms & gain \\
\midrule
q\_proj & $(1024, 8192)$ & 3.17 & 2.71 & 14.5\% \\
q\_proj & $(2048, 8192)$ & 1.82 & 1.62 & 11.2\% \\
up\_proj & $(1024, 8192)$ & 11.09 & 9.14 & 17.6\% \\
up\_proj & $(2048, 8192)$ & 6.44 & 5.62 & 12.8\% \\
\bottomrule
\end{tabular}
\end{center}

При этом численная точность не изменилась заметно: output MSE осталась на
уровне около $4.86 \times 10^{-6}$ для `q_proj` и $1.33 \times 10^{-6}$ для
`up_proj` на проверенных frontier-режимах. Лучший текущий режим остаётся
`(2048,8192)`, но теперь local path достигает `1.62` мс на `q_proj` и `5.62`
мс на `up_proj`, то есть находится уже примерно в `1.13x` и `1.05x` от fused
global full Triton соответственно.

Это наблюдение меняет интерпретацию ещё раз. Большая часть launch overhead уже
снята предыдущими grouped variants, и даже простой переход от `int32` к `int16`
для local indices даёт двузначный выигрыш. Следовательно, remaining gap теперь
почти целиком сосредоточен в том, как kernel обращается с palette values и как
организовано их staging/reuse внутри очень больших 2D blocks.

Последний проверенный nearby-probe касался уже самих local palette values.
После выигрыша от `int16`-индексов был проверен ещё один простой bandwidth move:
хранить локальный palette в `fp16` вместо `fp32`. Этот шаг тоже дал полезный
signal, но уже не столь однородный. Повторный замер на тех же frontier-режимах
показал следующую картину:

\begin{center}
\begin{tabular}{llrrr}
\toprule
Layer & $(group\_rows, group\_cols)$ & int16-only ms & fp16-palette ms & delta \\
\midrule
q\_proj & $(1024, 8192)$ & 2.71 & 2.48 & $-8.6\%$ \\
q\_proj & $(2048, 8192)$ & 1.62 & 1.57 & $-2.8\%$ \\
up\_proj & $(1024, 8192)$ & 9.14 & 8.44 & $-7.7\%$ \\
up\_proj & $(2048, 8192)$ & 5.62 & 5.89 & $+4.9\%$ \\
\bottomrule
\end{tabular}
\end{center}

Таким образом, `fp16`-palette помогает на трёх из четырёх проверенных режимов,
но слегка ухудшает самый большой `up\_proj` case. При этом численная точность
остается практически неизменной: output MSE остаётся порядка
$4.24 \times 10^{-6}$ для `q_proj` и $1.35 \times 10^{-6}$ для `up_proj`.

Это важное уточнение для системной интерпретации. Если переход на `int16`
indices был почти ``чистой победой'', то уменьшение bandwidth самих palette
values уже взаимодействует с shape-regime более сложно. Значит, финальный
remaining gap до fused global full path, по-видимому, нельзя снять ещё одним
простым dtype trick. Следующий наиболее правдоподобный источник выигрыша ---
явный palette staging/reuse внутри grouped local kernel.

Чтобы проверить, есть ли вообще у больших local palettes выраженный ``горячий''
префикс, был выполнен отдельный usage-profile experiment. Он показал, что на
frontier-группах `q_proj` действительно имеет сильную концентрацию: top-32
локальных значений покрывают около `90.8\%`, а top-64 около `95.7\%` всех
индексных обращений. Для `up_proj` распределение более плоское, но всё же
не равномерное: top-32 покрывают около `62.6\%`, а top-64 около `79.8\%`.

На основании этого был реализован прямой hot-prefix probe для local kernel.
Однако именно этот шаг дал ещё один важный отрицательный результат. Наивная
реализация через развернутую `where`-цепочку внутри Triton kernel оказалась
существенно хуже baseline grouped-local path. При `hot\_size=32` она была
примерно в `2.5x-3.5x` медленнее baseline. Уменьшение до `hot\_size=16`
примерно вдвое снизило этот overhead, но всё равно оставило kernel хуже
baseline на `21\%-57\%`.

Следовательно, usage-profile сам по себе не опровергает идею hot-prefix.
Напротив, он показывает, что особенно для `q_proj` горячий префикс действительно
существует. Но текущая форма эксплуатации этой структуры оказывается неверной.
Если продолжать эту ветку, нужен уже не branch-heavy substitution path, а
настоящий staged/shared hot-prefix design.

Был выполнен и последний ближайший discriminating check для этой ветки:
уменьшение префикса до `hot\_size=8`. Этот вариант действительно оказался лучше,
чем `hot\_size=16`, и почти достиг паритета на `q\_proj` в режиме
$(1024, 8192)$ (`2.486 ms` против `2.479 ms` у baseline grouped-local). Однако
ни на одном из проверенных frontier-режимов он всё равно не стал быстрее
baseline. Более того, повтор того же `hot\_size=8` probe на `2048` токенах не
спас идею, а закрыл её окончательно: naive hot-prefix kernel оказался лишь на
уровне примерно `0.59x-0.62x` от baseline grouped-local throughput. Тем самым
ветка ``наивный hot-prefix через branch-heavy substitution'' может считаться
закрытой.